\section{Abordagem}
O AHM é aplicado a um corpo cilíndrico, estendendo trabalhos anteriores de Guinovart-Sanjuan et al. (2015) e Guinovart-Sanjuan et al. (2017) ao aplicar o método descrito nesses trabalhos à situações comuns na mecânica, mais especificamente aos vasos de pressão. Utilizando o sistema de coordenadas cilíndricas ($r$, $\theta$, z) {conforme representado na Fig. \ref{fig:Cylindrical_coordinates}}, o comportamento elástico de estruturas heterogêneas com $N$ seções transversais cilíndricas concêntricas sob pressões internas e externas constantes é descrito. Para simplificar o problema, é considerado um estado de deformação plana (justificar a deformação plana), e as tensões de cisalhamento não são levadas em consideração.

A formulação matemática dos problemas locais e a derivação dos coeficientes efetivos são obtidas através do AHM. Os problemas locais são então resolvidos utilizando a linguagem de programação Python. Alguns resultados numéricos e uma análise da precisão da solução são realizados. Adicionalmente, uma simulação de compósito é conduzida implementando-se o método de elementos finitos (MEF) através do módulo FEniCS do Python. As fórmulas e resultados apresentados neste trabalho servem como pontos de referência valiosos para a validação dos resultados de cálculos experimentais e numéricos.



